% Generated draft: requires mathematical and source review.
\documentclass[11pt]{article}
\usepackage[margin=25mm]{geometry}
\usepackage{fontspec}
\setmainfont{Latin Modern Roman}
\usepackage{amsmath,amssymb,mathtools,unicode-math}
\setmathfont{Latin Modern Math}
\usepackage{enumitem,needspace,xurl,hyperref}
\setlength{\parindent}{0pt}
\setlength{\parskip}{5pt}
\setlength{\emergencystretch}{3em}
\newcommand{\N}{\mathbb{N}}
\newcommand{\Z}{\mathbb{Z}}
\newcommand{\Q}{\mathbb{Q}}
\newcommand{\R}{\mathbb{R}}
\newcommand{\C}{\mathbb{C}}
\newcommand{\F}{\mathbb{F}}
\newcommand{\A}{\mathbb{A}}
\newcommand{\PP}{\mathbb{P}}
\newcommand{\E}{\mathbb{E}}
\newcommand{\eps}{\varepsilon}
\newcommand{\dd}{\,\mathrm{d}}
\begin{document}
\section*{UnsolvedMath ID: 700088}
\textbf{Source problem code:} catalogue entry
\subsection{Definitions and mathematical statement}
A group homomorphism is a map between groups that preserves multiplication. A monomorphism of groups is an injective group homomorphism, and an epimorphism of groups is a surjective group homomorphism.

The source asks whether there exists a regular associative operation on the class of groups satisfying the weakened Mal'cev condition: monomorphisms of the factors of an arbitrary product can be glued together into a homomorphism of the whole product. The operation must fail the analogous condition for epimorphisms of the factors. The glued homomorphism in the stated monomorphism condition is not required to be injective.

Schematically, the requested existence statement is
\[
\exists\mathcal P\;\bigl(
\mathsf{RegularAssociative}(\mathcal P)
\land \mathsf{WeakMalcev}_{\mathrm{mono}}(\mathcal P)
\land \neg\mathsf{WeakMalcev}_{\mathrm{epi}}(\mathcal P)
\bigr).
\]
Here these predicates are placeholders for the source terminology, not independently specified axioms. The positive gluing condition concerns every arbitrary product and every admissible family of factor monomorphisms; failure of the analogous epimorphism condition requires a product and an admissible family of factor epimorphisms for which the corresponding gluing requirement fails.

\emph{Unresolved qualifications.} The supplied record does not define a regular associative operation, the product structure and its factor maps, or the compatibility and target requirements for gluing. It also does not specify the allowed index sets for arbitrary products or state the epimorphism analogue explicitly. Consequently, the displayed formula records the requested logical form, but a fully specified mathematical proposition requires the original definitions and conventions.
\subsection{Short English statement}
Does there exist a regular associative group product operation for which every admissible family of factor monomorphisms extends to a homomorphism of the product, while the analogous property for factor epimorphisms fails? The precise product and gluing conventions are missing from the supplied record.
\subsection{Sources}
\begin{itemize}
\item Local UnsolvedMath registry, supplied complete source record: permanent ID \(700088\), exact source problem code \texttt{catalogue entry}, title ``Kourovka Notebook Problem 2.9.'' The registry labels the problem open; this is supplied metadata, not an independently verified status assessment.
\item \emph{The Kourovka Notebook (Unsolved Problems in Group Theory)}, original problem 2.9, second issue (1967); proposed by O. N. Golovin. The supplied background identifies the twenty-first issue edition (2026), \texttt{arXiv:1401.0300v48}. The supplied collection description names E. I. Khukhro and V. D. Mazurov as editors.
\end{itemize}
\end{document}
